
\subsection{Execution-Based Feedback for Code Generation}

Execution feedback uses unit test outcomes as training signals. \textbf{CodeRL} \citep{le2022coderl} pioneered actor-critic reinforcement learning for code, using test pass/fail as episode-level rewards. A critic network predicts functional correctness, enabling critical sampling for code regeneration. CodeRL achieves 2.69 pass@1 on APPS with critic sampling.

\textbf{RLTF} \citep{liu2023rltf} extends execution feedback with multi-granularity signals. Beyond binary pass/fail, RLTF extracts fine-grained error locations from unit test output, providing line-level localization. Online RL with real-time data generation yields state-of-the-art results: 1.45 pass@1 on APPS without critic sampling. RLTF demonstrates that granularity matters---fine-grained feedback outperforms coarse rewards.

\textbf{PPOCoder} \citep{shojaee2023ppocoder} combines PPO-based RL with non-differentiable execution feedback and structure alignment. The framework is task-agnostic, achieving strong results across APPS and MBPP benchmarks.

These methods share a limitation: execution feedback provides no semantic guidance about \emph{why} code fails.

\subsection{AI-Generated Feedback for Code}

AI feedback uses LLM-generated critique for code improvement. \textbf{Self-Refine} \citep{madaan2023selfrefine} demonstrates zero-shot iterative refinement: the same LLM generates code, critiques it, and refines based on its own feedback. No training is required---refinement occurs at inference time. Self-Refine achieves +8.2\% improvement on code optimization tasks.

However, AI feedback suffers from low fidelity. Madaan et al.\ analyze Self-Refine failures: 33\% stem from identifying the wrong error location, and 61\% from proposing incorrect fixes.

\textbf{RefineCoder} \citep{zhou2025refinecoder} introduces Adaptive Critique Refinement using LLM-as-Judge and LLM-as-Critic for self-generated code refinement during training.

\subsection{The Gap}

No study directly compares execution and AI feedback under controlled conditions. Each method uses different base models, benchmarks, and training protocols. Furthermore, no prior work treats execution as a \emph{verification mechanism} for AI feedback.
